% 자동 생성 — scripts/build_tables.py. 손으로 고치지 않는다.
% 출처: output/results/12_spillover.json : 사다리_맞닿은동
\begin{table}[htbp]
  \centering
  \caption{같은 자료, 같은 창, 같은 사양. 기준점만 바꿨다 (맞닿은 동).}
  \label{tab:spill-bench}
  \begin{tabular}{lrrrr}
    \toprule
    \textbf{기준점} & \textbf{거래량 (\%)} & \textbf{$p$} & \textbf{가격 (\%)} & \textbf{$p$} \\
    \midrule
    1. 기준점 = 처치동 (선행연구의 시선) & 115.5 & $<0.001$ & 4.3 & 0.111 \\
    2. 기준점 = 강남·송파 비인접 동 & -6.8 & 0.331 & -3.6 & 0.055 \\
    3. 기준점 = 나머지 19개 구 & -9.0 & 0.088 & -5.0 & 0.001 \\
    4. 3 + 가격대 공통지지 & -19.0 & 0.021 & 0.3 & 0.748 \\
    \bottomrule
  \end{tabular}

  \begin{minipage}{\linewidth}\footnotesize
  \textit{주.} $+115.5\%$ 는 옆 동이 늘어서가 아니라 처치동이 $-57.8\%$ 로 줄어서다. 선행연구가 보고한 `인근 급등' 이 이것이다.
  \end{minipage}
\end{table}
